\documentclass[10pt]{article}

% Packages
\usepackage[a4paper,margin=1in]{geometry}
\usepackage{fancyhdr}
\usepackage{setspace}
\usepackage{graphicx}
\usepackage{float}
\usepackage{color}
\usepackage[
  style=apa,
  backend=biber
]{biblatex}
\addbibresource{master thesis.bib}

% Header setup
\pagestyle{fancy}
\fancyhf{} % clear default
\lhead{{\today}}
\rhead{Salome Bachmann}
\cfoot{\thepage}

% Optional: line under header
% \renewcommand{\headrulewidth}{0.4pt}

% Remove any automatic title page
\setlength{\parindent}{0pt}

\begin{document}

% If you want a title *inline*, not a title page:
{\centering \Large \textbf{Literature Review for Master Thesis}}\\[1pt]


\section{Avalanche Modelling}
\subsection{Numerical investigation of crack propagation regimes in snow fracture experiments \parencite{bobillierNumericalInvestigationCrack2024}}

\subsection{Micromechanical modeling of snow snow failure \parencite{bobillierMicromechanicalModelingSnow2020}}

\subsection{Modeling snow slab avalanches caused by weak-layer failure – Part 2: Coupled mixed-mode criterion for skier-triggered anticracks \parencite{rosendahlModelingSnowSlab2020}}

\subsection{Investigating the release and flow of snow avalanches at the slope-scale using a unified model based on the material point method \parencite{gaumeInvestigatingReleaseFlow2019}}

\subsection{Numerical investigation of the mixed-mode failure of snow \parencite{mulakNumericalInvestigationMixedmode2019}}

\subsection{Snow fracture in relation to slab avalanche release: critical state
for the onset of crack propagation \parencite{gaumeSnowFractureRelation2017a}}

\subsection{Modeling of crack propagation in weak snowpack layers using the discrete element method \parencite{gaumeModelingCrackPropagation2015}}

\subsection{A physical SNOWPACK model for the Swiss avalanche warning: Part I: numerical model \parencite{barteltPhysicalSNOWPACKModel2002}}
\textbf{Main Question of article:} How can the snowpack and its governing equations be modeled accurately but as simple as possible? \\
\textbf{Why this is relevant for the thesis:} Modeling snowpack evolution and structure is important to forecast avalanches (where are avalanches likely to form), but also to understand avalanche dynamics (in which layers do avalanches form) \\

\cite{barteltPhysicalSNOWPACKModel2002} describe that numerical modeling of the snowpack is important for avalanche prediction. The snowpack-model developed solves the partially governing equations using finite-element methods \parencite{barteltPhysicalSNOWPACKModel2002}. Understanding snowpack evolution and avalanche formation is also relevant for the thesis because the crack propagation which is modelled in the thesis depends on snow properties. According to \cite{barteltPhysicalSNOWPACKModel2002}, most avalanches occur during or after periods of intense snowfall. Snow settling gives an insight on mechanical stability of snow, which is also important to understand crack propagation: in stable layers cracks probably do not propagate as well as in layers which already have some instabilities. The microstructure and the layering of snow also needs to be modeled to detect discontinuities or layer boundaries, which is important to assess stability of layers \parencite{barteltPhysicalSNOWPACKModel2002}. \cite{barteltPhysicalSNOWPACKModel2002} also model the influence of solar radiation and warm winds on the energy exchange, the melting of snow and the subsequent transport of meltwater is important to understand the release of wet snow avalanches. \\

Existing snowpack models such as CROCUS rely on Trusdall's equation for mixture and postulate conservation of ice, water, water vapor and air of the snowpack, these models are simplified \parencite{barteltPhysicalSNOWPACKModel2002}. This means that the existing models for example do not describe settling of the snowpack, they assume the snowpack to be dry or the absence of snow metamorphism, which does not depcit the reality of the snowpack accurately \parencite{barteltPhysicalSNOWPACKModel2002}. An accurate and physical model of the snowpack is important for the thesis because modelled data should be comparable to field data measured by the DAS cable. This would not be the case if simplifications are made in the simulation of crack propagation in snow, the simulated data would look different from the measured data. The novel SNOWPACK model described by \cite{barteltPhysicalSNOWPACKModel2002} uses three constituents (ice, water and moist air) to describe the snow. These three all have the same temperature, which is the snow bulk temperature measured in meterological observations, assume all meltwater processes to be energy and mass conserving and assume that the weight of the snowpack is carried by a single ice lattice \parencite{barteltPhysicalSNOWPACKModel2002}. \\

\cite{barteltPhysicalSNOWPACKModel2002} use four governing equations to describe the snowpack: the energy transfer equation at the surface (conservation of mass of vapor and water and the bulk temperature), because all of the snow constituents have the same temperature, a momentum equation for settling of the snow because the weight of the snow is carried by a single ice lattice: The conservation of mass of ice is not considered because the numerical solution of the displacement is continuos because of the lagrangian coordinate system. Any change in the snow load results in a change of the stress state of the settlement equation. Snow cannot be assomed to be a plastic material because deformations caused by loading are irreversible, understanding this is important for crack formation and propagation. \cite{barteltPhysicalSNOWPACKModel2002} decsribe a microstructure-based viscosity formulation, plasticity theory is not applicable because of time-dependent yield surfaces and a directionality of plastic flow which is defined in relation to the yield surface. The initial conditions of the snowpack such as snow height are taken from meterological data, the energy exchange at the surface is the most important boundary condition \parencite{barteltPhysicalSNOWPACKModel2002}. Neumann boundary conditions are used if the temperature at the surface is close to the melting temperature, if the temperature is well below the melting temperature, Dirichlet boundary conditions are used \parencite{barteltPhysicalSNOWPACKModel2002}. \\

The melting and refreezing processes in the snowpack are essential in modelling snow structure changes, which have an effect on DAS signal propagation asn well as avalanche formation. \cite{barteltPhysicalSNOWPACKModel2002} model these by the numerical heat transfer solution. Because the snow temperature is assumed to be constant 0°C, if the heat transfer solution is greater than 0°C, the excess energy is assumed to be used for melting, if the heat transfer is smaller than 0°C, the energy is assumed to be used for refreezing \parencite{barteltPhysicalSNOWPACKModel2002}. \cite{barteltPhysicalSNOWPACKModel2002} used the snow structure measurements of winter 1999, which was a winter with numerous avalanches, to verify the SNOWPACK model developed.  Generally, there is good agreement between the measured and modelled data, which means that the model is physical, with some exceptions \parencite{barteltPhysicalSNOWPACKModel2002}. If the calculated snow height is different from the measured snow height, the model corrects this by adding elements to the finite element grid in case the calculated height is too low, or removing elements if snow hieght is overestimated \parencite{barteltPhysicalSNOWPACKModel2002}. This leads to mistakes if the snowpack settles quickly: if there is quick settlement, snow height decreases quickly because of densification, the model would remove elements at the surface without accounting for densification of the snow at the bottom \parencite{barteltPhysicalSNOWPACKModel2002}. Snow density and snow settlement have strong effects on heat conductivity of snow, \cite{barteltPhysicalSNOWPACKModel2002} mention that errors in either of these fields have strong effect on the entire simulation. \\

 The SNOWPACK model developed by \cite{barteltPhysicalSNOWPACKModel2002} includes modeling of the microstructure of snow as well as density discontinuities, which represent layer boundaries, which is important to understand snow structure as well as crack propagation related to avalaches. The multicomponent description developed allows for representation of mass and energy conservation \parencite{barteltPhysicalSNOWPACKModel2002}, which can be used to asess snow metamorphism related to temperature and moisture, but also to determine where wet snow avalanches form due to the presence of water. \cite{barteltPhysicalSNOWPACKModel2002} desrcibe the volumetric deformation of fresh snow by large strain description, the thermal boundary conditions which need to be used can be determined by meterological data. These data can also be used to determine if there is melting or refreezing of snow, which again is important for snow structure and stability. The laws for heat conductivity and viscosity related to microstructure show the relations between densification of snow and heat transfer \parencite{barteltPhysicalSNOWPACKModel2002}. 

\subsection{A numerical model to simulate snow-cover stratigraphy for operational avalanche forecasting \parencite{brunNumericalModelSimulate1992}}
\textbf{Main Question of article:} How do snow conditions such as temperature gradient and water content influence snow metamporphism and how can this be implemented in the CROCUS-snowpack model? \\
\textbf{Why this is relevant for the thesis:} Snow metamorphism changes grain properties, which is important for the formation of weak layers which can trigger avalanche release, but also overall stability of the snowpack. Furthermore, layers of different grain properties will depict DAS signals differently. \\ 

When a snowpack is subjected to a change, like a change in temperature or pressure conditions, the snow grains change \parencite{brunNumericalModelSimulate1992}. Snow metamorphism changes the grain sizes and shapes, which changes mechanical stability of the snow \parencite{brunNumericalModelSimulate1992}. This is relevant for avalanche modelling because certain types of snow metamorphism can weaken layers in snow, such layers can act as failure planes and make release of avalanches easier. Snow metamorphism is affected by the water content of the snow, the capillary pressure of the water between the snow grains, the radius of curvature of the grains (shape), as well as the temperature gradient \parencite{brunNumericalModelSimulate1992}. \\

\cite{brunNumericalModelSimulate1992} describes the main limitation of past methamorphism models as assuming the snowpack to have homogenous shapes and sizes, which are often spherical or simplified geometric shapes, this does not represent snow accurately however, especially fresh snow is very heterogenous. \cite{brunNumericalModelSimulate1992} conducted experiments for different snow types and conditions to determine new equations to describe metamorphism. These experiments showed that the rate and type of metamorphism does not depend on crystal type but rather on the temperature gradient \parencite{brunNumericalModelSimulate1992}. A gradient below 5°C/m results in rounded crystals, a gradient above 5°C/m results in facteted crystals, which transform to depth hoar, if a gradient of more than 15°C/m is applied \parencite{brunNumericalModelSimulate1992}. This is relevant for avalanche modelling, because it is critical to understand how snow grains evolve to asess their stability, which can be used to determine where weak layers will form. These experiments were conducted for dry and wet snow conditions.\\

\cite{brunNumericalModelSimulate1992} describe a new description of snow, which takes into account the dendricity, sphericity and size of the grains. When snow falls, new layers with different properties are added to the snowpack \parencite{brunNumericalModelSimulate1992}. Simulating these different layers is again important to assess the stability as well as interactions of the different layers. \cite{brunNumericalModelSimulate1992} combined the observations made in the laboratory with the CROCUS snowpack model to simulate these different layers based on meterological observations. Using MEPRA (a system for avalanche risk forecasting), the shear stress of the different layers was estimated at a 45° angle. This is important for avalanche forcasting, because the strength of different layers can be assessed and it can be determined under which stress the layers fail. 

\subsection{A discrete numerical model for granular assemblies \parencite{cundallDiscreteNumericalModel1979}}
\textbf{Main Question of article:} How can the forces and the resulting displacements be modeled in a granular medium?\\
\textbf{Why this is relevant for the thesis:} Snow consists of grains, forces at the surface such as mechanical loading can cause displacement of the individual grains. If this happens at a very large scale, it is an avalanche. \\

The distinct element method described by \cite{cundallDiscreteNumericalModel1979} is a numerical modeling method to model the mechanical behavior of a material composed of disks & spheres. This is relevant to the thesis because snow is composed of individual snow grains, which interact with each other at their contacts, all the grains together form the snow pack. It is important to model the behavior and displacement of the grains to model avalanche dynamics because avalanches form through loading or failure in a weak layer, which then propagates through the entire snowpack \parencite{gaumeSnowFractureRelation2017a}. \\

 Since a snowpack is a medium composed by individual snow grains, which are not entriely compacted (contain air between the grains) and interact with each other at grain contacts,  the medium described by \cite{cundallDiscreteNumericalModel1979} represents snow well. It is a granular medium composed of distinct particles, which interact with each other at their contacts and form a porous medium \parencite{cundallDiscreteNumericalModel1979}. The loading and unloading of such materials is complex and stresses cannot be measured easily because of the interactions of different grains \parencite{cundallDiscreteNumericalModel1979}. \\

The distinct elememt method can model particles of any shape \parencite{cundallDiscreteNumericalModel1979}, which is important for snow modelling because the snow grains are not always spherical. \cite{cundallDiscreteNumericalModel1979} model the transient interaction of particles by calculating the contact forces and displacements of the stressed discs (particles) caused by disturbances at the boundaries. In the avalanche case discussed in the thesis, the distrubance is represented by loading of the surface, which can be caused by a skier, which causes failure in a weak layer, or cracking in the snowpack. This results in stress propagation through the snowpack, which results in displacement of the particles, which is the avalanche motion. The resultant foces on the disks are calculated only by the discs which are in direct contact (all the disks that touch one disk cause its resultant force) \parencite{cundallDiscreteNumericalModel1979}. \cite{cundallDiscreteNumericalModel1979} describe a model only for dry materials, this means this is only applicable for dry snow avalanches which contain no water between the snow grains. \\

Numerical modelling, such as described by \cite{cundallDiscreteNumericalModel1979}is advantageus over lab tests and experiments because it is less time consuming and any data can be accessed at any stage, which means that parameters such as loading, particle size, and properties such as stiffness and shape can be changed at any stage. In the avalanche case, this means that simulations can be run for different snow properties and loading scenarios, for example a simulation for older, more compacted snow can be run as well as a simulation for fresh snow. The numerical model for spheres is based on methods described by Serrano and Roderiguez-Ortiz (1973, I CANNOT FIND THIS PAPER ONLINE). The forces at the interfaces between the spheres are calculated by incremental displacements of the centers of particles, where shape changes of the individual particles are neglible \parencite{cundallDiscreteNumericalModel1979}. \textcolor{magenta}{How can the center of the grain be determined if the grain has an irrgular shape?} The matrix representing contact stiffness between the grains described by \cite{cundallDiscreteNumericalModel1979} needs to be recalculated whenever a new contact is made or broken, meaning when there is a crack in the snowpack or when new contacts between snow grains are made. The calculations of forces are based on Newtons 2nd law, which describes the movement of a particle as the result of the forces acting on it. The deformation is neglected because the deformation of individual particles is small compared to the deformation of the medium as a whole \parencite{cundallDiscreteNumericalModel1979}. This means that the deformation of the individual snow grains is small compared to the deformation of the entire snowpack which represents the avalanche movement. \textcolor{magenta}{What about the fusing / compaction of grains which changes snow properties?} The relative displacement of a particle is calculated by the sum of the force increments acting on it, partciles can also overlap, in this case the contact foces and the interaction forces also need to be calculated \parencite{cundallDiscreteNumericalModel1979}. The location of the rigid boundaries (walls) described by \cite{cundallDiscreteNumericalModel1979} needs to be specified before modelling. This could be a possible limitation for the avalanche case because there is only one rigid "wall" which is the interface of snow and the rock below it. 



\section{Crack Propagation}

\section{Relevant Source Mechanisms}


\newpage
\section{Bibliography} 
\typeout{}
\printbibliography


\end{document}